\documentclass[a4paper,12pt]{article}
\usepackage{fontspec} % Пакет для работы со шрифтами в XeLaTeX
\setmainfont{Liberation Serif} % Теперь эта команда сработает
\usepackage[russian]{babel}

% Математические пакеты
\usepackage{amsmath, amsfonts, amssymb}
\usepackage{geometry}
\geometry{left=3cm, right=1.5cm, top=2cm, bottom=2cm}
\usepackage{geometry}
\usepackage{graphicx}
\usepackage{tabularx}
\usepackage{booktabs}
\usepackage{indentfirst}

\geometry{top=2cm, bottom=2cm, left=2.5cm, right=1.5cm}

\begin{document}

\thispagestyle{empty}

% Шапка: Название университета и факультет
\begin{center}
    {\fontsize{14}{16}\selectfont \textbf{Университет ИТМО}} \\
    {\fontsize{12}{14}\selectfont Физико-технический мегафакультет} \\
    {\fontsize{12}{14}\selectfont Физический факультет}
\end{center}

% Место для логотипа (справа сверху)
\begin{flushright}
    \vspace{-3cm}
    \Huge \textbf{ИТМО} % Плейсхолдер логотипа
\end{flushright}

\vspace{1.5cm}

% Блок с информацией о студентах и проверке
\noindent
\begin{tabular}{ll}
    Поток \makebox[5cm]{ФИЗ-3 Э БИТ 1.2.1} & К работе допущен\underline{\makebox[3cm]{}} \\[0.3cm]
    Студенты: {\makebox[5.5cm]{Пахомов Владимир Юрьевич,}} & Работа выполнена\underline{\makebox[3cm]{}} \\
    Сивков иван Александрович & \\
    Преподаватель: Терещенко Георгий Викторович & Отчет принят\underline{\makebox[3.5cm]{}}
\end{tabular}

\vspace{2cm}

% Заголовок работы
\begin{center}
    {\fontsize{18}{22}\selectfont \textbf{Рабочий протокол и отчет}} \\[0.5cm]
    {\fontsize{18}{22}\selectfont \textbf{по лабораторной работе № 4.11}} \\[1.5cm]
    
    {\fontsize{16}{20}\selectfont \textbf{Определение основных характеристик дифракционной}} \\
    {\fontsize{16}{20}\selectfont \textbf{решетки}}
\end{center}




% Дальнейшее содержание отчета (Цели, Задачи и т.д.)
\section{Цель работы}
1) Изучение характеристик дифракционной решетки.

\section{Задачи, решаемые при выполнении работы}
1) Экспериментальное определение угловой дисперсии решетки; \\
2) Экспериментальное определение разрешающей способности решетки.

\section{Объект исследования}
Дифракционная решётка.

\section{Метод экспериментального исследования}
Метод основан на измерении угловых положений спектральных линий, полученных с помощью дифракционной решётки, для определения ее основных характеристик. При прохождении света через решётку происходит дифракция и интерференция волн, в результате чего излучение различных длин волн распространяется под разными углами. Измеряя углы дифракции для известных длин волн, определяют период решётки. На основе полученных данных рассчитываются угловая дисперсия и разрешающая способность решетки, а также число штрихов на единицу длины.








\section{Рабочие формулы и исходные данные}

\begin{equation}
d \sin(\varphi) = m\lambda \tag{1}
\end{equation}
— условие возникновения интерференционных максимумов для дифракционной решетки.

\begin{equation}
n = \frac{1}{d} \tag{2}
\end{equation}
— число штрихов, нанесенных на 1 мм ширины решетки.

\begin{equation}
D = \frac{d\varphi}{d\lambda} = \frac{\Delta\varphi}{\Delta\lambda} \tag{3}
\end{equation}
— угловая дисперсия (экспериментальная).

\begin{equation}
D = \frac{d\varphi}{d\lambda} = \frac{m}{d \cos(\varphi)} \tag{4}
\end{equation}
— угловая дисперсия (теоретическая).

\begin{equation}
R = \frac{\lambda}{\delta\lambda} \tag{5}
\end{equation}
— разрешающая способность (определение).

\begin{equation}
R = mN \tag{6}
\end{equation}
— разрешающая способность дифракционной решетки.

\begin{equation}
\varphi = \frac{N_2 - N_1}{2} \tag{7}
\end{equation}
— угол дифракции.

\vspace{0.5cm}
\textbf{Исходные данные для длин волн (ртутная лампа):}
\begin{itemize}
    \item Зеленая линия: $\lambda_{зел} = 546,1$ нм
    \item Синяя линия: $\lambda_{син} = 435,8$ нм
\end{itemize}

\section{Результаты прямых и косвенных измерений и их обработки}

\begin{center}
\begin{tabular}{|c|c|c|c|c|c|}
\hline
 & $\varphi_1$ & $\varphi_2$ & $\varphi_3$ & $\varphi_{ср}$ & $m$ \\ \hline
$N_1^{зел}$ & $18^\circ 26'$ & $19^\circ 17'$ & $18^\circ 57'$ & $18^\circ 53'$ & 1 \\ \hline
$N_1^{син}$ & $15^\circ 32'$ & $15^\circ 15'$ & $15^\circ 28'$ & $15^\circ 25'$ & 1 \\ \hline
$N_2^{зел}$ & $339^\circ 28'$ & $340^\circ 32'$ & $339^\circ 42'$ & $339^\circ 54'$ & -1 \\ \hline
$N_2^{син}$ & $344^\circ 20'$ & $344^\circ 40'$ & $344^\circ 34'$ & $344^\circ 31'$ & -1 \\ \hline
\end{tabular}

\medskip
Таблица 1 - Результаты измерений.
\end{center}

\noindent Рассчитаем углы дифракции по формуле (7):
\begin{equation*}
\varphi_{зел} = \frac{18^\circ 53' - (339^\circ 54' - 360^\circ)}{2} = \frac{18^\circ 53' - (-20^\circ 06')}{2} = 19^\circ 30'
\end{equation*}
\begin{equation*}
\varphi_{син} = \frac{15^\circ 25' - (344^\circ 31' - 360^\circ)}{2} = \frac{15^\circ 25' - (-15^\circ 29')}{2} = 15^\circ 27'
\end{equation*}

\noindent Рассчитаем период решетки с помощью формулы (1) для зеленой линии ($\lambda = 546,1$ нм):
\begin{equation*}
d = \frac{m\lambda}{\sin(\varphi_{зел})} = \frac{546,1}{\sin(19^\circ 30')} = \frac{546,1}{0,3338} \approx 1636,0 \text{ нм}
\end{equation*}

\noindent Рассчитаем число штрихов на 1 мм с помощью формулы (2):
\begin{equation*}
n = \frac{1}{d} = \frac{1}{1636,0 \cdot 10^{-6} \text{ мм}} \approx 611,2 \text{ мм}^{-1}
\end{equation*}

\noindent Рассчитаем угловую дисперсию решетки по экспериментальным данным (формула 3):
\begin{equation*}
\Delta\varphi = \varphi_{зел} - \varphi_{син} = 19^\circ 30' - 15^\circ 27' = 4^\circ 03' = 243'
\end{equation*}
\begin{equation*}
\Delta\lambda = \lambda_{зел} - \lambda_{син} = 546,1 - 435,8 = 110,3 \text{ нм}
\end{equation*}
\begin{equation*}
D = \frac{243'}{110,3 \text{ нм}} \approx 2,20 \frac{'}{\text{нм}}
\end{equation*}

\noindent Рассчитаем угловую дисперсию решетки по теоретической формуле (4):
\begin{equation*}
D = \frac{m}{d \cos(\varphi)} = \frac{1}{1636,0 \cdot \cos(19^\circ 30')} = \frac{1}{1636,0 \cdot 0,9426} \approx 6,48 \cdot 10^{-4} \frac{\text{рад}}{\text{нм}} \approx 2,23 \frac{'}{\text{нм}}
\end{equation*}
Результаты вычислений хорошо согласуются.

\noindent Рассчитаем разрешающую способность решетки по формуле (6), приняв ширину решетки $L = 300$ мм:
\begin{equation*}
N = n \cdot L = 611,2 \cdot 300 = 183360 \text{ (полное число штрихов)}
\end{equation*}
\begin{equation*}
R = mN = 1 \cdot 183360 = 183360
\end{equation*}




\section{Расчет погрешностей}

\noindent Относительная погрешность определения периода решетки $d$ находится путем дифференцирования выражения $d = \frac{m\lambda}{\sin(\varphi)}$:
\begin{equation*}
\ln(d) = \ln(m\lambda) - \ln(\sin\varphi)
\end{equation*}
\begin{equation*}
\frac{\Delta d}{d} = \left| \frac{\partial}{\partial \varphi} \ln(\sin\varphi) \right| \cdot \Delta \varphi = \frac{\cos\varphi}{\sin\varphi} \Delta \varphi = \frac{\Delta \varphi}{\tan\varphi}
\end{equation*}

\noindent Подставим значения для зеленой линии ($\varphi = 19^\circ 30'$, $d = 1636,0$ нм):
\begin{itemize}
    \item $\tan(19^\circ 30') \approx 0,3541$
    \item $\Delta \varphi = 1' \approx 2,91 \cdot 10^{-4}$ рад
\end{itemize}

\noindent Рассчитаем относительную погрешность:
\begin{equation*}
\varepsilon = \frac{\Delta d}{d} = \frac{2,91 \cdot 10^{-4}}{0,3541} \approx 0,000822 \approx 0,082\%
\end{equation*}

\noindent Рассчитаем абсолютную погрешность периода решетки:
\begin{equation*}
\Delta d = d \cdot \varepsilon = 1636,0 \cdot 0,000822 \approx 1,35 \text{ нм}
\end{equation*}

\noindent Переведем значения в миллиметры для итоговой записи (как в примере):
\begin{itemize}
    \item $d = 1636,0 \text{ нм} = 0,0016360 \text{ мм}$
    \item $\Delta d = 1,35 \text{ нм} \approx 0,0000014 \text{ мм}$
\end{itemize}

\noindent \textbf{Итоговый результат:}
\begin{equation*}
d = (0,0016360 \pm 0,0000014) \text{ мм}
\end{equation*}




\section{Вывод}

\noindent \textbf{Вывод:} В ходе лабораторной работы были изучены основные характеристики дифракционной решетки. Экспериментально измерены угловые положения спектральных линий ртутной лампы (зеленой и синей), на основе которых определен период решетки $d \approx 1636,0$ нм и число штрихов на единицу длины $n \approx 611,2 \text{ мм}^{-1}$.

\medskip
В результате расчетов была определена угловая дисперсия решетки двумя способами:
\begin{enumerate}
    \item По экспериментальным данным: $D_{эксп} \approx 2,20 \frac{'}{\text{нм}}$;
    \item По теоретической формуле: $D_{теор} \approx 2,23 \frac{'}{\text{нм}}$.
\end{enumerate}
Полученные значения хорошо согласуются между собой (относительное расхождение составляет менее 1.5\%), что подтверждает корректность проведенных измерений и расчетов.

\medskip
Также была рассчитана разрешающая способность решетки ($R = 183360$ для первого порядка спектра при ширине $L = 300$ мм), что позволяет оценить высокую способность данной оптической системы разделять близкие по длине волны спектральные линии.

\medskip
Таким образом, результаты эксперимента подтверждают теоретические закономерности дифракции и интерференции света на дифракционной решетке.

\vfill
\begin{flushright}
    4
\end{flushright}


\end{document}